\section{Mode F4 -- Mode manuel}
\label{secModeF4}

\subsection{Conditions d'entrée}

Le mode F4 est atteint depuis le mode A1, sur demande de marche manuelle
formulée par l'opérateur (exclusive de la demande de marche automatique
menant au mode F1).

\subsection{Comportement attendu}

Ce mode est destiné aux besoins de maintenance : il permet à l'opérateur de
commander directement, depuis le pupitre opérateur, le déplacement de
chacun des axes du robot, simultanément ou non, dans le respect des
butées articulaires de chaque axe.

\begin{itemize}
    \item \textbf{Convoyeur / guichet} : l'accès au guichet est interdit
    pendant toute la durée du mode manuel.
    \item \textbf{Robot} : les commandes de mouvement issues du pupitre
    opérateur ne sont prises en compte que si le guichet est vide. Une
    demande de mouvement contradictoire (par exemple une demande
    simultanée de monter et de descendre un même axe) ne doit provoquer
    aucun mouvement.
    \item \textbf{Comportement pendant un mouvement en cours} : si
    l'opérateur relâche la commande à l'origine du mouvement (ou formule
    une nouvelle demande contradictoire), le mouvement en cours doit être
    interrompu et l'installation doit se tenir prête à accepter une
    nouvelle demande de mouvement valide. Si l'opérateur maintient sa
    commande jusqu'à la fin de course demandée, le robot atteint la
    position limite correspondante avant de se remettre en attente d'une
    nouvelle demande.
    \item \textbf{IHM} : la situation de mode manuel doit être signalée en
    permanence par une signalisation dédiée, clignotante.
\end{itemize}

\subsection{Conditions de sortie}

\begin{itemize}
    \item Vers le mode \textbf{A6} (\emph{Mise en conditions initiales}),
    sur demande de retour aux conditions initiales formulée par
    l'opérateur, à condition qu'aucun mouvement ne soit alors en cours.
    \item Vers le mode \textbf{D1}, si une défaillance de communication est
    détectée.
\end{itemize}
